% =====================================================================
%  Chapitre 7 : Gestion des requetes
% =====================================================================
\chapter{Gérer les requêtes et les vues}
\label{ch:requetes}

\chapitreresume{Une vue reçoit une requête et renvoie une réponse. Dans ce
chapitre, vous explorez l'objet \code{request}, les codes de statut, les
raccourcis \code{redirect} et \code{get\_object\_or\_404}, les convertisseurs
d'URL, la recherche, les filtres et la pagination. Le catalogue de
\gesibiblio{} devient une vraie application qui interroge la base de données.}

\begin{objectifs}
\begin{itemize}
  \item Lire les informations d'une requête HTTP dans l'objet \code{request}.
  \item Choisir le bon type de réponse et le bon code de statut.
  \item Utiliser \code{get\_object\_or\_404}, \code{redirect} et
        \code{reverse}.
  \item Maîtriser les convertisseurs d'URL et les paramètres nommés.
  \item Construire une liste avec recherche, filtre, tri et pagination.
  \item Protéger une vue selon la méthode HTTP.
  \item Écrire les pages d'erreur 404 et 500.
  \item Découvrir les vues basées sur les classes en alternative.
\end{itemize}
\end{objectifs}

\section{Le trajet complet d'une requête}

Au chapitre précédent, les vues renvoyaient des données écrites à la main. Elles
vont maintenant interroger la base. Le trajet devient complet :

\begin{enumerate}
  \item le navigateur envoie une requête \code{GET /livres/} ;
  \item les middlewares préparent \code{request} ;
  \item le routeur choisit la vue \code{livre\_liste} ;
  \item la vue interroge les modèles avec l'ORM ;
  \item la base renvoie les lignes ;
  \item la vue transmet un contexte au gabarit ;
  \item le gabarit produit du HTML ;
  \item Django renvoie une \code{HttpResponse} au navigateur.
\end{enumerate}

\begin{retenir}
Une vue est le chef d'orchestre. Elle ne contient ni HTML ni SQL : elle valide la
requête, interroge les modèles, choisit un gabarit et renvoie une réponse. Cette
séparation rend le code testable et lisible.
\end{retenir}

\section{L'objet request}

Django construit un objet \code{HttpRequest} pour chaque requête. Les attributs
les plus utiles :

\begin{table}[htbp]
\centering
\small
\begin{tabularx}{\textwidth}{L{4.6cm} X}
\toprule
\textbf{Attribut ou méthode} & \textbf{Contenu} \\
\midrule
\code{request.method} & Méthode HTTP : \code{"GET"}, \code{"POST"}, etc. \\
\code{request.GET} & Paramètres de l'URL, après le point d'interrogation. \\
\code{request.POST} & Champs du corps, pour un formulaire \code{POST}. \\
\code{request.FILES} & Fichiers téléversés. \\
\code{request.path} & Chemin demandé, par exemple \code{/livres/}. \\
\code{request.GET.urlencode} & Paramètres encodés pour construire un lien. \\
\code{request.headers} & En-têtes HTTP, insensibles à la casse. \\
\code{request.COOKIES} & Cookies envoyés par le navigateur. \\
\code{request.user} & Utilisateur connecté, ou \code{AnonymousUser}. \\
\code{request.session} & Données de session. \\
\code{request.body} & Corps brut de la requête, en octets. \\
\code{request.content\_type} & Type de contenu annoncé. \\
\code{request.is\_secure()} & Vrai si la requête arrive en HTTPS. \\
\bottomrule
\end{tabularx}
\caption{Les attributs principaux de l'objet request.}
\end{table}

\subsection{GET et POST}

\begin{codebox}[python]{Lire les paramètres}
def recherche(request):
    mot = request.GET.get("q", "").strip()
    page = request.GET.get("page", "1")
    return render(
        request,
        "catalogue/recherche.html",
        {"mot": mot, "page": page},
    )
\end{codebox}

\begin{description}[leftmargin=4.3cm,style=nextline]
  \item[\code{request.GET["q"]}] Lève \code{KeyError} si le paramètre manque.
  \item[\code{request.GET.get("q")}] Renvoie \code{None} si le paramètre manque.
  \item[\code{request.GET.get("q", "")}] Renvoie une chaîne vide par défaut :
  c'est la forme la plus sûre pour une recherche.
  \item[\code{.strip()}] Supprime les espaces autour de la valeur.
\end{description}

\begin{attention}
Ne faites jamais confiance aux valeurs de \code{request.GET} ou
\code{request.POST}. Ce sont des chaînes envoyées par le client. Un utilisateur
peut modifier l'URL, envoyer \code{page=abc} ou \code{id=-5}. Validez le type, la
longueur et les valeurs autorisées avant d'utiliser une donnée.
\end{attention}

\subsection{Les en-têtes et la négociation de contenu}

\begin{codebox}[python]{Lire un en-tête}
navigateur = request.headers.get("User-Agent", "inconnu")
langue = request.headers.get("Accept-Language", "fr")
souhaite_json = "application/json" in request.headers.get("Accept", "")
\end{codebox}

\section{Construire une réponse}

Django propose plusieurs classes de réponse. Elles héritent toutes de
\code{HttpResponse}.

\begin{table}[htbp]
\centering
\small
\begin{tabularx}{\textwidth}{L{5.1cm} X}
\toprule
\textbf{Classe ou fonction} & \textbf{Usage} \\
\midrule
\code{HttpResponse} & Réponse HTML ou texte brut. \\
\code{JsonResponse} & Sérialise un dictionnaire en JSON. \\
\code{HttpResponseRedirect} & Redirection \code{302}. \\
\code{HttpResponsePermanentRedirect} & Redirection \code{301}. \\
\code{HttpResponseNotFound} & Erreur \code{404}. \\
\code{HttpResponseForbidden} & Erreur \code{403}. \\
\code{FileResponse} & Envoie un fichier en flux. \\
\code{StreamingHttpResponse} & Envoie un flux progressif. \\
\code{redirect} & Raccourci de redirection. \\
\code{render} & Charge un gabarit et renvoie une réponse. \\
\bottomrule
\end{tabularx}
\caption{Les réponses les plus utilisées.}
\end{table}

\subsection{Le code de statut}

\begin{codebox}[python]{Choisir un statut}
from django.http import HttpResponse, JsonResponse


def sante(request):
    return JsonResponse({"statut": "ok", "version": "1.0"})


def page_absente(request):
    return HttpResponse(
        "<h1>Contenu introuvable</h1>",
        status=404,
    )


def methode_refusee(request):
    return HttpResponse(
        "Cette action n'accepte que POST.",
        status=405,
    )
\end{codebox}

\begin{bonne}
Un code de statut correct aide les navigateurs, les moteurs de recherche et les
outils de supervision. Une page \og introuvable \fg{} doit renvoyer \code{404},
pas \code{200} avec un joli message. Une création réussie renvoie souvent
\code{302} vers la fiche, ou \code{201} dans une API.
\end{bonne}

\subsection{Rediriger}

Après un formulaire validé, on redirige vers une autre URL. Cela évite qu'un
rafraîchissement ne renvoie le formulaire.

\begin{codebox}[python]{Redirection}
from django.shortcuts import redirect, render
from .models import Livre


def livre_creation_rapide(request):
    if request.method == "POST":
        titre = request.POST.get("titre", "").strip()
        if titre:
            livre = Livre.objects.create(
                titre=titre,
                slug=titre.lower().replace(" ", "-"),
                isbn="9780000000000",
                auteur_id=1,
                annee_publication=2024,
            )
            return redirect(livre)
    return render(request, "catalogue/livres/creation.html")
\end{codebox}

\begin{description}[leftmargin=4.4cm,style=nextline]
  \item[\code{redirect("catalogue:livre\_liste")}] Redirige vers une route nommée.
  \item[\code{redirect(livre)}] Utilise \code{livre.get\_absolute\_url()}.
  \item[\code{redirect("/livres/")}] Redirige vers un chemin littéral, à éviter.
\end{description}

\begin{note}
Le motif \og POST puis redirection puis GET \fg{} s'appelle le patron
\emph{Post/Redirect/Get}. Il empêche le navigateur de proposer de renvoyer le
formulaire lors d'un rafraîchissement. Nous l'utiliserons systématiquement avec
les formulaires.
\end{note}

\section{Les objets et les 404}

\begin{codebox}[python]{Récupérer un objet}
from django.shortcuts import get_object_or_404, render
from .models import Livre


def livre_detail(request, slug):
    livre = get_object_or_404(
        Livre.objects.select_related("auteur").prefetch_related("categories"),
        slug=slug,
    )
    return render(
        request,
        "catalogue/livres/detail.html",
        {"livre": livre},
    )
\end{codebox}

\code{get\_object\_or\_404} exécute la requête et lève une \code{Http404} si aucun
objet ne correspond. Django renvoie alors la page 404, ou la page de débogage en
développement. Sans ce raccourci, il faudrait écrire :

\begin{codebox}[python]{Équivalent sans raccourci}
from django.http import Http404

try:
    livre = Livre.objects.get(slug=slug)
except Livre.DoesNotExist:
    raise Http404("Ce livre n'existe pas.")
\end{codebox}

\begin{bonne}
Utilisez \code{get\_object\_or\_404} pour toute fiche accessible par URL. Cela
évite une erreur 500 lorsqu'un identifiant n'existe pas et donne au visiteur une
réponse propre. Il existe aussi \code{get\_list\_or\_404} pour les listes.
\end{bonne}

\begin{attention}
\code{get()} sans filtre unique peut lever \code{MultipleObjectsReturned}. Un slug
ou un ISBN sont uniques : ce sont de bons paramètres d'URL. N'utilisez pas
\code{titre} comme identifiant, deux livres peuvent avoir le même titre.
\end{attention}

\section{Les URL et leurs convertisseurs}

Une URL peut capturer une partie du chemin et la passer à la vue.

\begin{codebox}[python]{catalogue/urls.py}
from django.urls import path
from . import views

app_name = "catalogue"

urlpatterns = [
    path("", views.accueil, name="accueil"),
    path("a-propos/", views.a_propos, name="a_propos"),
    path("livres/", views.livre_liste, name="livre_liste"),
    path("livres/<slug:slug>/", views.livre_detail, name="livre_detail"),
    path(
        "categories/<slug:slug>/",
        views.categorie_detail,
        name="categorie_detail",
    ),
    path("auteurs/<int:pk>/", views.auteur_detail, name="auteur_detail"),
]
\end{codebox}

\begin{table}[htbp]
\centering
\small
\begin{tabularx}{\textwidth}{L{3.1cm} L{4.6cm} X}
\toprule
\textbf{Convertisseur} & \textbf{Exemple} & \textbf{Valeur reçue} \\
\midrule
\code{str} & \code{<str:nom>} & Toute chaîne sans barre oblique. \\
\code{int} & \code{<int:pk>} & Un entier positif ou nul. \\
\code{slug} & \code{<slug:slug>} & Lettres, chiffres, tirets et soulignés. \\
\code{uuid} & \code{<uuid:identifiant>} & Un UUID sous forme de chaîne. \\
\code{path} & \code{<path:chemin>} & Une chaîne qui peut contenir des barres. \\
\bottomrule
\end{tabularx}
\caption{Convertisseurs d'URL fournis par Django.}
\end{table}

\begin{codebox}[python]{Vue avec paramètre entier}
def auteur_detail(request, pk):
    auteur = get_object_or_404(
        Auteur.objects.prefetch_related("livres"),
        pk=pk,
    )
    return render(
        request,
        "catalogue/auteurs/detail.html",
        {"auteur": auteur},
    )
\end{codebox}

\begin{note}
Le convertisseur \code{int} refuse les valeurs négatives et les lettres : Django
renvoie directement une 404, la vue n'est jamais appelée. C'est une première
barrière de validation, gratuite et fiable.
\end{note}

\subsection{Les paramètres nommés}

\begin{codebox}[python]{Paramètres nommés dans une URL}
path(
    "livres/<int:annee>/<slug:slug>/",
    views.livre_par_annee,
    name="livre_par_annee",
),
\end{codebox}

La vue doit alors accepter les deux arguments :

\begin{codebox}[python]{Signature correspondante}
def livre_par_annee(request, annee, slug):
    ...
\end{codebox}

\begin{attention}
L'ordre des \code{path} compte. Une route trop générale placée avant une route
précise capture les URL. Par exemple, \code{<str:nom>/} placé avant
\code{livres/} empêcherait d'atteindre le catalogue. Placez les routes
spécifiques en premier.
\end{attention}

\subsection{reverse et resolve}

\begin{codebox}[python]{Construire et analyser une URL}
from django.urls import resolve, reverse

url = reverse("catalogue:livre_detail", kwargs={"slug": "le-petit-prince"})
# /livres/le-petit-prince/

vue = resolve("/livres/le-petit-prince/")
vue.func          # la fonction livre_detail
vue.kwargs        # {'slug': 'le-petit-prince'}
vue.url_name      # 'livre_detail'
vue.namespace     # 'catalogue'
\end{codebox}

\begin{bonne}
Dans le code Python, utilisez \code{reverse}. Dans un gabarit, utilisez
\code{\{\% url \%\}}. Ces deux fonctions construisent les liens à partir du nom
de route et évitent les chemins codés en dur.
\end{bonne}

\section{La liste des livres}

Nous allons construire une liste complète : recherche, filtre par catégorie, tri et
pagination.

\subsection{La vue}

\begin{codebox}[python]{catalogue/views.py, extrait}
from django.core.paginator import Paginator
from django.db.models import Count, Q
from django.shortcuts import get_object_or_404, render

from .models import Auteur, Categorie, Livre

TRI_AUTORISE = {
    "titre": "titre",
    "recent": "-annee_publication",
    "auteur": "auteur__nom",
}


def livre_liste(request):
    livres = (
        Livre.objects.select_related("auteur")
        .prefetch_related("categories")
    )

    mot = request.GET.get("q", "").strip()
    if mot:
        livres = livres.filter(
            Q(titre__icontains=mot)
            | Q(auteur__nom__icontains=mot)
            | Q(isbn__icontains=mot)
        )

    categorie_slug = request.GET.get("categorie", "").strip()
    if categorie_slug:
        livres = livres.filter(categories__slug=categorie_slug)

    tri = request.GET.get("tri", "titre")
    livres = livres.order_by(TRI_AUTORISE.get(tri, "titre"))

    paginator = Paginator(livres, 6)
    page = paginator.get_page(request.GET.get("page"))

    categories = Categorie.objects.annotate(
        nb_livres=Count("livres")
    )

    contexte = {
        "page_obj": page,
        "livres": page.object_list,
        "categories": categories,
        "mot": mot,
        "categorie_active": categorie_slug,
        "tri": tri,
        "nombre_resultats": paginator.count,
    }
    return render(request, "catalogue/livres/liste.html", contexte)
\end{codebox}

\begin{attention}
Ne confondez pas deux usages du symbole \code{|}. À l'intérieur d'un
\code{filter}, \code{Q(...) | Q(...)} combine des conditions dans la même
requête : c'est la forme recommandée. Entre deux QuerySets,
\code{queryset1 | queryset2} fusionne deux résultats, ne dédoublonne pas toujours
et empêche parfois d'enchaîner un filtre. Pour une recherche sur plusieurs champs,
utilisez \code{Q}.
\end{attention}

\subsection{Comprendre la pagination}

\begin{codebox}[python]{Le Paginator}
from django.core.paginator import Paginator

paginator = Paginator(Livre.objects.all(), 6)
page = paginator.get_page(3)
\end{codebox}

\begin{description}[leftmargin=4.6cm,style=nextline]
  \item[\code{paginator.count}] Nombre total d'objets.
  \item[\code{paginator.num\_pages}] Nombre total de pages.
  \item[\code{paginator.page\_range}] Plage de numéros de page.
  \item[\code{page.number}] Numéro de la page courante.
  \item[\code{page.object\_list}] Objets de la page.
  \item[\code{page.has\_previous()}] Vrai s'il existe une page avant.
  \item[\code{page.has\_next()}] Vrai s'il existe une page après.
  \item[\code{page.previous\_page\_number}] Numéro de la page précédente.
  \item[\code{page.next\_page\_number}] Numéro de la page suivante.
  \item[\code{paginator.get\_page(n)}] Renvoie la page \code{n}, ou la première ou
  la dernière page si la valeur est invalide. Ne lève pas d'erreur : c'est la
  méthode à utiliser avec un paramètre utilisateur.
\end{description}

\begin{bonne}
\code{get\_page()} accepte \code{None}, une chaîne vide, \code{"abc"} ou un numéro
hors limites sans lever d'exception. Il renvoie la page appropriée. Préférez-le à
\code{paginator.page()}, qui lève \code{EmptyPage} ou \code{PageNotAnInteger}.
\end{bonne}

\subsection{Le gabarit de liste}

\begin{codebox}[html]{templates/catalogue/livres/liste.html}
{% extends "base.html" %}
{% load biblio %}

{% block titre %}Catalogue | GesiBiblio{% endblock %}

{% block contenu %}
  <h1>Catalogue</h1>

  <form method="get" class="formulaire filtre">
    <div class="champ">
      <label for="q">Rechercher</label>
      <input type="search" id="q" name="q" value="{{ mot }}"
             placeholder="Titre, auteur ou ISBN">
    </div>
    <div class="champ">
      <label for="categorie">Catégorie</label>
      <select id="categorie" name="categorie">
        <option value="">Toutes</option>
        {% for categorie in categories %}
          <option value="{{ categorie.slug }}"
            {% if categorie.slug == categorie_active %}selected{% endif %}>
            {{ categorie.nom }} ({{ categorie.nb_livres }})
          </option>
        {% endfor %}
      </select>
    </div>
    <div class="champ">
      <label for="tri">Trier par</label>
      <select id="tri" name="tri">
        <option value="titre" {% if tri == "titre" %}selected{% endif %}>
          Titre
        </option>
        <option value="recent" {% if tri == "recent" %}selected{% endif %}>
          Année décroissante
        </option>
        <option value="auteur" {% if tri == "auteur" %}selected{% endif %}>
          Auteur
        </option>
      </select>
    </div>
    <button class="bouton bouton-principal" type="submit">Filtrer</button>
  </form>

  <p>
    {{ nombre_resultats }} livre{{ nombre_resultats|pluralize }}
    trouvé{{ nombre_resultats|pluralize }}.
  </p>

  <div class="grille-livres">
    {% for livre in livres %}
      {% carte_livre livre %}
    {% empty %}
      <p>Aucun livre ne correspond à votre recherche.</p>
    {% endfor %}
  </div>

  {% if page_obj.has_other_pages %}
    <nav class="pagination" aria-label="Pagination du catalogue">
      {% if page_obj.has_previous %}
        <a href="{% url_page page_obj.previous_page_number %}">Précédent</a>
      {% endif %}
      <span>
        Page {{ page_obj.number }} sur {{ page_obj.paginator.num_pages }}
      </span>
      {% if page_obj.has_next %}
        <a href="{% url_page page_obj.next_page_number %}">Suivant</a>
      {% endif %}
    </nav>
  {% endif %}
{% endblock %}
\end{codebox}

\begin{note}
La balise \code{url\_page} a été créée au chapitre précédent. Elle conserve les
paramètres \code{q}, \code{categorie} et \code{tri} en changeant seulement
\code{page}. Sans elle, la page suivante perdrait les filtres.
\end{note}

\subsection{La fiche d'un livre}

\begin{codebox}[html]{templates/catalogue/livres/detail.html}
{% extends "base.html" %}

{% block titre %}{{ livre.titre }} | GesiBiblio{% endblock %}

{% block contenu %}
  <article class="fiche-livre">
    {% if livre.couverture %}
      <img src="{{ livre.couverture.url }}"
           alt="Couverture de {{ livre.titre }}">
    {% else %}
      <div class="couverture-vide" aria-hidden="true">
        {{ livre.titre|first }}
      </div>
    {% endif %}

    <div>
      <p class="surtitre">{{ livre.get_statut_display }}</p>
      <h1>{{ livre.titre }}</h1>
      {% if livre.sous_titre %}
        <p class="auteur">{{ livre.sous_titre }}</p>
      {% endif %}

      <dl>
        <dt>Auteur</dt>
        <dd>
          <a href="{% url 'catalogue:auteur_detail' livre.auteur.pk %}">
            {{ livre.auteur }}
          </a>
        </dd>
        <dt>Éditeur</dt>
        <dd>{{ livre.editeur|default:"Non renseigné" }}</dd>
        <dt>Année</dt>
        <dd>{{ livre.annee_publication }}</dd>
        <dt>ISBN</dt>
        <dd>{{ livre.isbn }}</dd>
        <dt>Catégories</dt>
        <dd>
          {% for categorie in livre.categories.all %}
            <a href="{% url 'catalogue:categorie_detail' categorie.slug %}">
              {{ categorie.nom }}
            </a>{% if not forloop.last %}, {% endif %}
          {% empty %}
            Aucune catégorie
          {% endfor %}
        </dd>
        <dt>Exemplaires</dt>
        <dd>
          {{ livre.exemplaires_disponibles }} disponible{{
          livre.exemplaires_disponibles|pluralize }} sur
          {{ livre.nombre_exemplaires }}
        </dd>
      </dl>

      {% if livre.resume %}
        <h2>Résumé</h2>
        <div>{{ livre.resume|linebreaks }}</div>
      {% endif %}
    </div>
  </article>
{% endblock %}
\end{codebox}

\begin{attention}
Sur la ligne des exemplaires, le gabarit coupe une expression sur deux lignes. En
HTML, un retour à la ligne entre deux mots affiche un espace : c'est correct. Mais
ne coupez jamais une balise \code{\{\{ ... \}\}} au milieu de son contenu, cela
provoque une erreur de gabarit.
\end{attention}

\subsection{Les vues de catégorie et d'auteur}

\begin{codebox}[python]{catalogue/views.py, extrait}
def categorie_detail(request, slug):
    categorie = get_object_or_404(Categorie, slug=slug)
    livres = (
        categorie.livres.select_related("auteur")
        .prefetch_related("categories")
    )
    return render(
        request,
        "catalogue/categories/detail.html",
        {"categorie": categorie, "livres": livres},
    )


def auteur_detail(request, pk):
    auteur = get_object_or_404(Auteur, pk=pk)
    livres = auteur.livres.select_related("auteur")
    return render(
        request,
        "catalogue/auteurs/detail.html",
        {"auteur": auteur, "livres": livres},
    )
\end{codebox}

\begin{note}
\code{categorie.livres} est un gestionnaire lié par \code{related\_name}. Il se
comporte comme \code{Livre.objects}, mais limité aux livres de la catégorie. On
peut donc enchaîner \code{filter}, \code{select\_related} et \code{order\_by}.
\end{note}

\section{Protéger une vue selon la méthode HTTP}

Certaines actions ne doivent accepter que \code{POST}. Django fournit des
décorateurs.

\begin{codebox}[python]{Décorateurs de méthode}
from django.views.decorators.http import (
    require_GET,
    require_POST,
    require_http_methods,
)


@require_GET
def livre_liste(request):
    ...


@require_POST
def livre_supprimer(request, pk):
    ...


@require_http_methods(["GET", "POST"])
def livre_modifier(request, pk):
    ...
\end{codebox}

Si la méthode ne correspond pas, Django renvoie \code{405 Method Not Allowed}.
C'est plus propre qu'un message maison et plus sûr qu'une suppression déclenchée
par un simple lien.

\begin{attention}
Ne supprimez jamais une donnée avec un lien \code{<a href="...">}. Un moteur de
recherche ou un antivirus peut suivre le lien et supprimer vos enregistrements. La
suppression se fait avec un formulaire \code{POST}, une confirmation et, idéalement,
une vérification de permission.
\end{attention}

\section{Les pages d'erreur}

En développement, Django affiche une page de débogage détaillée. En production,
avec \code{DEBUG = False}, il utilise des gabarits.

\begin{codebox}[html]{templates/404.html}
{% extends "base.html" %}

{% block titre %}Page introuvable | GesiBiblio{% endblock %}

{% block contenu %}
  <section class="hero">
    <div>
      <p class="surtitre">Erreur 404</p>
      <h1>Cette page est introuvable.</h1>
      <p>Le livre que vous cherchez a peut-être changé de rayon.</p>
      <a class="bouton bouton-principal"
         href="{% url 'catalogue:accueil' %}">Retour à l'accueil</a>
    </div>
  </section>
{% endblock %}
\end{codebox}

\begin{codebox}[html]{templates/500.html}
{% extends "base.html" %}

{% block titre %}Erreur serveur | GesiBiblio{% endblock %}

{% block contenu %}
  <section class="hero">
    <div>
      <p class="surtitre">Erreur 500</p>
      <h1>Une erreur est survenue.</h1>
      <p>Réessayez dans quelques instants. L'incident a été signalé.</p>
      <a class="bouton" href="/">Retour à l'accueil</a>
    </div>
  </section>
{% endblock %}
\end{codebox}

\begin{attention}
Le gabarit \code{500.html} ne doit pas dépendre de données ou de balises
personnalisées. Si une erreur se produit, le rendu de la page d'erreur peut lui
aussi échouer. Restez simple et n'utilisez pas \code{url} dans \code{500.html} :
préférez un chemin littéral ou une page statique.
\end{attention}

Pour tester, créez un fichier \code{.env} ou modifiez temporairement
\code{settings.py} :

\begin{codebox}[python]{gesibiblio/settings.py, extrait}
DEBUG = False
ALLOWED_HOSTS = ["127.0.0.1", "localhost"]
\end{codebox}

Puis lancez \code{python manage.py runserver --insecure} pour servir les fichiers
statiques malgré \code{DEBUG = False}. N'oubliez pas de remettre
\code{DEBUG = True} ensuite.

\section{Les vues basées sur les classes}

Les vues basées sur les fonctions sont le fil principal de ce module. Django
propose aussi des classes génériques, plus concises pour les cas standard.

\begin{codebox}[python]{catalogue/views.py, alternative CBV}
from django.views.generic import DetailView, ListView

from .models import Livre


class LivreListView(ListView):
    model = Livre
    template_name = "catalogue/livres/liste.html"
    context_object_name = "livres"
    paginate_by = 6

    def get_queryset(self):
        return Livre.objects.select_related("auteur")


class LivreDetailView(DetailView):
    model = Livre
    template_name = "catalogue/livres/detail.html"
    context_object_name = "livre"
\end{codebox}

\begin{codebox}[python]{catalogue/urls.py, alternative CBV}
path("livres/", views.LivreListView.as_view(), name="livre_liste"),
path("livres/<slug:slug>/", views.LivreDetailView.as_view(), name="livre_detail"),
\end{codebox}

\begin{description}[leftmargin=4.4cm,style=nextline]
  \item[\code{ListView}] Récupère une liste, la pagine et la place dans le
  contexte. Le nom par défaut est \code{object\_list} ; ici, il devient
  \code{livres} grâce à \code{context\_object\_name}.
  \item[\code{DetailView}] Récupère un objet à partir de \code{pk} ou
  \code{slug}. Le nom par défaut est \code{object}.
  \item[\code{get\_queryset()}] Point d'entrée pour filtrer ou optimiser la
  requête.
\end{description}

\begin{bonne}
Pour un débutant, les vues basées sur les fonctions sont plus faciles à lire et à
déboguer. Les classes deviennent intéressantes quand vous répétez le même schéma
de liste ou de détail. Ne mélangez pas les deux styles dans une même vue.
\end{bonne}

\begin{piege}
Si vous remplacez une vue fonction par une classe, pensez à \code{.as\_view()} dans
\code{urls.py}. Sans cela, Django reçoit une classe et non une fonction : vous
obtenez une erreur \code{TypeError} ou un comportement inattendu.
\end{piege}

\section{Exercices}

\begin{exercice}[ : fiche catégorie]
Écrivez la vue \code{categorie\_detail} et son gabarit. La page affiche le nom de
la catégorie, sa description et la liste de ses livres. Ajoutez un compteur et un
lien vers chaque fiche.
\end{exercice}

\begin{correction}
\begin{codebox}[html]{templates/catalogue/categories/detail.html}
{% extends "base.html" %}
{% load biblio %}

{% block titre %}{{ categorie.nom }} | GesiBiblio{% endblock %}

{% block contenu %}
  <h1>{{ categorie.nom }}</h1>
  {% if categorie.description %}
    <p>{{ categorie.description }}</p>
  {% endif %}
  <p>
    {{ livres|length }} livre{{ livres|length|pluralize }} dans cette catégorie.
  </p>
  <div class="grille-livres">
    {% for livre in livres %}
      {% carte_livre livre %}
    {% empty %}
      <p>Aucun livre dans cette catégorie.</p>
    {% endfor %}
  </div>
{% endblock %}
\end{codebox}

La vue est donnée plus haut. Vérifiez la route
\code{categories/<slug:slug>/} et le \code{related\_name} du modèle.
\end{correction}

\begin{exercice}[ : recherche par année]
Ajoutez un filtre \code{annee} à la liste des livres. L'utilisateur saisit une
année et la liste ne montre que les livres publiés cette année-là. Si la valeur
n'est pas un nombre, ignorez le filtre.
\end{exercice}

\begin{correction}
\begin{codebox}[python]{catalogue/views.py, extrait}
annee = request.GET.get("annee", "").strip()
if annee.isdigit():
    livres = livres.filter(annee_publication=int(annee))
\end{codebox}

\begin{codebox}[html]{templates/catalogue/livres/liste.html, extrait}
<div class="champ">
  <label for="annee">Année</label>
  <input type="number" id="annee" name="annee" value="{{ request.GET.annee }}">
</div>
\end{codebox}

\code{isdigit()} évite une erreur \code{ValueError} si l'utilisateur envoie
\code{annee=abc}. Le filtre est simplement ignoré.
\end{correction}

\begin{exercice}[ : réponse JSON]
Créez une vue \code{/api/livres/} qui renvoie la liste des titres et des auteurs
au format JSON. Utilisez \code{JsonResponse} et \code{values()}.
\end{exercice}

\begin{correction}
\begin{codebox}[python]{catalogue/views.py, extrait}
from django.http import JsonResponse


def api_livres(request):
    donnees = list(
        Livre.objects.values("titre", "auteur__nom", "annee_publication")
    )
    return JsonResponse({"livres": donnees})
\end{codebox}

\begin{codebox}[python]{catalogue/urls.py, extrait}
path("api/livres/", views.api_livres, name="api_livres"),
\end{codebox}

\code{values()} renvoie des dictionnaires, directement sérialisables.
\code{JsonResponse} ajoute l'en-tête \code{Content-Type: application/json}.
\end{correction}

\begin{retenir}
\begin{itemize}
  \item \code{request} contient la méthode, les paramètres, les fichiers, les
        en-têtes et l'utilisateur.
  \item \code{get\_object\_or\_404} transforme un objet absent en page 404.
  \item \code{redirect} évite le renvoi de formulaire au rafraîchissement.
  \item Les convertisseurs \code{int} et \code{slug} valident l'URL avant la vue.
  \item \code{Paginator} et \code{get\_page} paginent sans erreur.
  \item Une action qui modifie des données passe par \code{POST}.
  \item Les vues basées sur les classes sont une alternative, pas une obligation.
\end{itemize}
\end{retenir}
